\documentclass[11pt]{article}
\usepackage{parskip}
\usepackage{color}

\title{Hamantashen}
\author{Nate Zuk}

\begin{document}
\maketitle

\textcolor{green}{(3/15/2018)}

Based on the dough recipe from \href{https://www.chabad.org/library/article_cdo/aid/1366/jewish/Traditional-Hamantaschen.htm}{Chabad's website}.

\textbf{Ingredients:}

\begin{itemize}
\item Flour = ~5 cups
\item Eggs = 4
\item Sugar = 1 cup
\item Vegetable oil = 1/2 cup
\item Lemon juice = ~1/2 cup \textcolor{red}{(It was a big lemon)}
\item Vanilla = 1 tsp
\item Baking powder = 2 tsp \textcolor{red}{(I accidentally use 1 tsp of baking soda and 1 tsp of baking powder, but that turned out fine)}
\end{itemize}

Preheat the oven to 350 F.

Beat together eggs and sugar in a large bowl.  Add oil, lemon juice, and vanilla and mix together.  Add flour 1 cup at a time, mixing after each addition.  After about 3-4 cups, when the dough starts to maintain shape, add the baking powder and mix well.  Once the dough holds together and is not so sticky, turn it out on a floured surface and knead briefly to get the dough to hold together.

Separate the dough into two halves.  For each half, roll out the dough on a floured surface until about 1/8 to 1/4 inch in width.  Rotate the dough as you roll it out to keep in from sticking to the surface (but carefully, because the tough can come apart pretty easily).  Use a glass to cut the dough into ~3 inch diameter circles.  Add the filling to the center (about 1/2 Tbsp), pinch the top of the dough, then pinch the bottom ends to form a triangle.  If the filling is a bit runny, bend the dough up to keep the filling in as you pinch.

Bake the cookies on a cookie sheet with parchment paper for 20 minutes.  Let cool on the sheet before removing (~4 min).

\end{document}